In this work, we develop a SCFT framework that explicitly incorporates the sliding degree of freedom $f$, and apply it to L--S TT copolymers to study their microphase separation behavior.
By treating the grafting position as a continuously adjustable variable and using Gauss--Legendre quadrature to sample the topological space, this framework captures the contributions of multiple sliding configurations.
On this basis, we construct complete SCFT phase diagrams for three representative side-chain fractions, $\phi_{\mathrm a}=0.25$, $0.50$, and $0.75$, showing how the sliding mechanism reshapes phase boundaries at the macroscopic scale.
Near the spinodal, the preferred topology $f_\Lambda^\ast$ obtained from SCFT is consistent with the linear-response prediction of RPA, confirming the validity of the topology-selection picture.
Beyond this weak-segregation limit, however, SCFT retains the spatially resolved conditional topology distribution $P_{\mathcal R(\bm r)}(f)$ inside ordered phases, enabling the evaluation of the sliding-distribution entropy loss and revealing how sliding topology responds to local interfacial geometry and heterogeneous filling environments.



The RPA spinodal difference between the freely sliding and most unstable fixed-topology references provides an effective conformational entropy scale, $\Delta S_{\mathrm{conf}}$, which varies nonmonotonically with $\phi_{\mathrm a}$ and reaches its maximum near the crossover between star and linear topology locking.
In the ordered phase, the normalized local topology distributions quantify the sliding distribution entropy loss, $\Delta S_{\mathrm{Slide}}$, relative to uniform sliding.
Together, these measures connect the conformational response near the homogeneous instability with the redistribution of the sliding ensemble after microphase separation, revealing two coupled consequences of the same sliding degree of freedom.



Furthermore, within ordered phases, the sliding degree of freedom enables local topological redistribution, allowing L--S TT copolymers to respond to variations in interfacial curvature and accommodate heterogeneous local filling environments.
In the $\mathrm{O}^{70}$ phase, the topology redistribution varies with the local interfacial environment, with distinct responses in the A-core and A-shell geometries and a dependence of the preferred topology on local curvature.
In spherical phases, spatially varying preferred topologies coexist with broad local topology distributions, giving rise to topology homogenization while allowing the sliding degree of freedom to accommodate heterogeneous local packing environments.
With increasing $N\chi$, the stabilizing effect of sliding continues to strengthen in BCC but becomes strongest at intermediate segregation and weakens at higher $N\chi$ in FCC, indicating that its effectiveness depends on how broadly topology regulation can operate within the ordered structure.

More generally, the topology response evolves from a global preference at weak segregation toward increasingly local regulation as segregation strengthens, reflecting the coupled conformational and interfacial response.

This evolution allows sliding topology to accommodate local interfacial and packing constraints under stronger segregation.


Overall, this work reveals, at the field-theoretic level, how the sliding topological degree of freedom in L--S TT copolymers links thermodynamic stability, interfacial geometry, and local packing constraints.
The results show that sliding topology is not only an additional degree of freedom in the entropy--enthalpy competition of microphase separation, but also a mechanism for stabilizing complex morphologies, regulating local interfacial properties, and programming soft-matter structures.
In this context, the local distribution $P_{\mathcal R(\bm r)}(f)$ and the regional statistic $P_{\Omega}(f)$ provide spatially resolved measures for characterizing topology responses to curvature and local filling environments associated with packing frustration, offering a theoretical basis for the molecular design of tunable-topology copolymers. 
More broadly, the sliding degree of freedom may also provide a useful screening variable for the design of conventional fixed-topology systems.
By mapping the locally preferred topology $f_{\mathcal R}^{\ast}$ together with the strength of topology selection quantified by the local topology distribution, the present framework may help identify a reduced set of fixed molecular architectures, or combinations of architectures, that are better matched to heterogeneous local environments.

The framework developed here can be extended to more complex spherical phases, such as Frank--Kasper phases, where heterogeneous packing constraints across inequivalent local environments are central to phase stability.
It may also be applied to curvature-controlled nanostructures and to the phase behavior of slide-ring polymers, mechanically interlocked polymers, and other dynamic-topology copolymer systems, providing a general field-theoretic route for designing polymeric materials and programmable soft-matter structures.
